\documentclass[a4paper,11pt]{article}
\usepackage[margin=24mm]{geometry}
\usepackage{amsmath,amssymb,bm}
\usepackage{mathtools}
\usepackage{hyperref}
\usepackage{enumitem}
\usepackage{booktabs}
\usepackage{physics}
\setlength{\parskip}{0.4em}
\setlength{\parindent}{0em}

\title{OLED Cavity仕様書\\
\large verifyTM記法に統一したコード再監査版}
\author{}
\date{\today}

\begin{document}
\maketitle

\section{目的と今回の修正点}
本書は \texttt{phase2p7.tar} と \texttt{verify.tar} の実コードを照合し，
記法を \texttt{verifyTM.tex} の流儀（B/T面定義式 + rewrap写像式）に統一した仕様書である．

主修正は以下：
\begin{itemize}[leftmargin=1.5em]
\item Bottom / Top の「巻き戻し層」と「終端反射の定義面」を厳密に分離して記述
\item 比較式
\[
r_{\mathrm{rewrap,bot/top}}^{(\sigma)} \stackrel{?}{=} r_{\mathrm{full,bot/top}}^{(\sigma)}
\]
に統一
\item 全モジュール再点検結果（記述ミス有無・実装上注意）を反映
\end{itemize}

\section{記法（verifyTM流儀）}
偏光 \(\sigma\in\{s,p\}\)，
\[
u=\frac{k_\parallel}{k_0},\qquad k_0=\frac{2\pi}{\lambda}.
\]

PyMoosh整合規約（\(\mu=1,\ \varepsilon=n^2\)）:
\[
k_{z,j}=k_0\sqrt{\varepsilon_j-u^2},\qquad \Im(k_{z,j})\ge 0,
\]
\[
b_j^{(s)}=k_{z,j},\qquad b_j^{(p)}=\frac{k_{z,j}}{\varepsilon_j},
\]
\[
r_{j,k}^{(\sigma)}=\frac{b_j^{(\sigma)}-b_k^{(\sigma)}}{b_j^{(\sigma)}+b_k^{(\sigma)}}.
\]

再帰演算子 \(\mathcal W\) を
\[
\mathcal W\!\left(n_{\mathrm{inc}};\{(n_1,d_1),\dots,(n_N,d_N)\};r_{\mathrm{term}}^{(\sigma)}\right)
\]
と書く．内部再帰は
\[
r^{(N)}=r_{\mathrm{term}}^{(\sigma)},
\]
\[
r^{(j-1)}
=\frac{r_{j-1,j}^{(\sigma)}+r^{(j)}e^{2ik_{z,j}d_j}}
{1+r_{j-1,j}^{(\sigma)}r^{(j)}e^{2ik_{z,j}d_j}},
\quad (j=N,\dots,1)
\]
で，最終値 \(r^{(0)}\) を返す．

\section{B/T/S面の厳密定義}
\begin{itemize}[leftmargin=1.5em]
\item B面（Bottom終端反射の定義面）: \(pHTL|ITO\)（pHTL側から見た反射）
\item T面（Top終端反射の定義面）: \(ETL|cathode1\)（ETL側から見た反射）
\item S面: EML中心
\end{itemize}

注：「B=ITO/pHTL」「T=ETL/cathode1」は境界名の略記であり，
反射係数の入射側は上記のとおり pHTL側 / ETL側で定義する．

\section{Bottom/Topのrewrapとfull-stack比較式}

\subsection{Bottom}
終端反射:
\[
r_B^{(\sigma)}(u,\lambda)
=
r^{(\sigma)}\!\left(
\mathrm{incident}=pHTL,\;
\mathrm{layers}=[ITO,anode],\;
\mathrm{substrate}=substrate
\right).
\]

rewrap（EML|EBL面へ）:
\[
r_{\mathrm{rewrap,bot}}^{(\sigma)}(u,\lambda)
=
\mathcal W\!\left(
EML;\;
[(EBL,d_{EBL}),(HTL,d_{HTL}),(Rprime,d_{Rprime}),(pHTL,d_{pHTL})];\;
r_B^{(\sigma)}
\right).
\]

full-stack:
\[
r_{\mathrm{full,bot}}^{(\sigma)}(u,\lambda)
=
r^{(\sigma)}\!\left(
\mathrm{incident}=EML,\;
\mathrm{layers}=[EBL,HTL,Rprime,pHTL,ITO,anode],\;
\mathrm{substrate}=substrate
\right).
\]

検証条件:
\[
r_{\mathrm{rewrap,bot}}^{(\sigma)}
\stackrel{?}{=}
r_{\mathrm{full,bot}}^{(\sigma)}.
\]

\subsection{Top}
終端反射:
\[
r_T^{(\sigma)}(u,\lambda)
=
r^{(\sigma)}\!\left(
\mathrm{incident}=ETL,\;
\mathrm{layers}=[cathode1,cathode,CPL],\;
\mathrm{substrate}=air
\right).
\]

rewrap（EML|ETL面へ）:
\[
r_{\mathrm{rewrap,top}}^{(\sigma)}(u,\lambda)
=
\mathcal W\!\left(
EML;\;
[(ETL,d_{ETL})];\;
r_T^{(\sigma)}
\right).
\]

full-stack:
\[
r_{\mathrm{full,top}}^{(\sigma)}(u,\lambda)
=
r^{(\sigma)}\!\left(
\mathrm{incident}=EML,\;
\mathrm{layers}=[ETL,cathode1,cathode,CPL],\;
\mathrm{substrate}=air
\right).
\]

検証条件:
\[
r_{\mathrm{rewrap,top}}^{(\sigma)}
\stackrel{?}{=}
r_{\mathrm{full,top}}^{(\sigma)}.
\]

\section{S面（EML中心）への写像とFP proxy}
EML境界反射を \(r_{EML/EBL}^{(\sigma)}, r_{EML/ETL}^{(\sigma)}\) とすると
\[
\Delta z=\frac{d_{EML}}{2},\qquad
k_{z,EML}=k_0\sqrt{\varepsilon_{EML}-u^2},
\]
\[
r_u^{(S,\sigma)}=r_{EML/EBL}^{(\sigma)}e^{2ik_{z,EML}\Delta z},\qquad
r_d^{(S,\sigma)}=r_{EML/ETL}^{(\sigma)}e^{2ik_{z,EML}\Delta z},
\]
\[
F^{(\sigma)}(\lambda,u)=
\frac{|1+r_u^{(S,\sigma)}|^2}{|1-r_u^{(S,\sigma)}r_d^{(S,\sigma)}|^2}.
\]

\section{モジュール別仕様（再監査反映）}

\subsection{(1) \texttt{tmm\_rewrap\_utils.py}}
\begin{itemize}[leftmargin=1.5em]
\item 上記 \(\mathcal W\) を実装
\item TE/TM は \(b=k_z,\;k_z/\varepsilon\) 規約
\item \(k_z\) 分岐は \(\Im(k_z)\ge 0\) で固定
\item Bottom再帰層: \([EBL,HTL,Rprime,pHTL]\)
\item Top再帰層: \([ETL]\)
\end{itemize}
入力: \texttt{n0,d,lam\_nm,u,rB\_pol,rT\_pol,pol}，
出力: \((ru_S,rd_S)\) および \texttt{F}．

\subsection{(2) \texttt{oled\_cavity**}}
対象:
\begin{itemize}[leftmargin=1.5em]
\item \texttt{oled\_cavity\_pymoosh\_sourceplane\_FP\_finetune.py}
\item \texttt{oled\_cavity\_pymoosh\_sourceplane\_FP\_finetune\_compat\_TMMrewrap.py}
\end{itemize}

位相条件:
\[
L_{\mathrm{target}}=m\frac{\lambda_0}{2}
-\frac{\lambda_0}{4\pi}(\phi_b+\phi_e),\qquad \phi_b=\pi,
\]
\[
s=\frac{L_{\mathrm{target}}-L_{\mathrm{fixed}}}{L_{\mathrm{unit}}}.
\]
B面位相固定は \(r_B\leftarrow-|r_B|\) で実施．
ETL粗探索（1 nm）\(\to\)細探索（0.1 nm），
目的関数は EML中心と \(|E|^2\)ピークの距離．

実装差：
\begin{itemize}[leftmargin=1.5em]
\item \texttt{finetune.py}: B/T参照の位相シフト式（\texttt{F\_BT}と理論等価）
\item \texttt{compat\_TMMrewrap.py}: 再帰TMM写像後に \(F\) 評価
\end{itemize}

監査メモ：\texttt{compat\_TMMrewrap.py} は fine-scan で early return（互換挙動）を保持．

\subsection{(3) \texttt{reproduce**}}
対象:
\texttt{reproduce\_phase2\_step2\_u\_dependence\_pymoosh\_sourceplane\_FP\_TMMrewrap.py}

\[
F_{TE/TM}(\lambda,u)\ \text{をS面再帰TMMで算出し，}
\quad K_{\mathrm{iso}}=\frac{2}{3}K_h+\frac{1}{3}K_v.
\]

Fig.4変換:
\[
u_{\mathrm{req}}=
\frac{\lambda[\mathrm{nm}]\times 10^{-9}}{2\pi}k_\parallel.
\]

実装上，
\(\theta\) は \texttt{theta\_from\_u\_branch\_fixed} で枝固定．
既定では \(K_h=F_{TE},\ K_v=10^{-10}\)（簡略版）．

\subsection{(4) \texttt{verify}（bottom/top）}
対象:
\begin{itemize}[leftmargin=1.5em]
\item \texttt{tmm\_rewrap\_utils\_verify.py}
\item \texttt{test\_verify\_bottom\_rewrap\_vs\_fullstack\_v3.py}
\item \texttt{test\_verify\_top\_rewrap\_vs\_fullstack\_v3.py}
\end{itemize}

誤差定義:
\[
\Delta r=|r_{\mathrm{rewrap}}-r_{\mathrm{full}}|.
\]

実測（同梱スクリプト実行）:
\begin{itemize}[leftmargin=1.5em]
\item Bottom: \(\max|\Delta r|_{TE}=3.333\times10^{-15}\), \(\max|\Delta r|_{TM}=1.355\times10^{-15}\)
\item Top: \(\max|\Delta r|_{TE}=1.746\times10^{-14}\), \(\max|\Delta r|_{TM}=4.220\times10^{-14}\)
\end{itemize}

\section{再監査結論}
\begin{enumerate}[leftmargin=1.8em]
\item Bottom/Top巻き戻しは verify と整合：
Bottomは終端 \(r_B\) 生成後に \([EBL,HTL,Rprime,pHTL]\) を写像，
Topは終端 \(r_T\) 生成後に \([ETL]\) を写像．
\item 前回問題は実装ミスではなく仕様書側の省略表現（誤読誘発）．
\item 実装注意点として，
\texttt{compat\_TMMrewrap.py} の early return と
\texttt{reproduce**} の \(K_h/K_v\) 簡略化を明記すべき．
\end{enumerate}

\section{実行コマンド}
\begin{verbatim}
# verify
python test_verify_bottom_rewrap_vs_fullstack_v3.py
python test_verify_top_rewrap_vs_fullstack_v3.py

# phase2p7
python oled_cavity_pymoosh_sourceplane_FP_finetune.py
python oled_cavity_pymoosh_sourceplane_FP_finetune_compat_TMMrewrap.py
python reproduce_phase2_step2_u_dependence_pymoosh_sourceplane_FP_TMMrewrap.py
\end{verbatim}

\end{document}
